\scanpage{105}
\mspage{3.4}
unramified over $S$. Since $\overline D=\overline R$, we have a natural epimorphism $D\otimes_R S\xrightarrow{\varphi}\overline S$. Let $\underline M=\operatorname{Ker}\varphi$ and set $\widetilde S=(D\otimes_R S)_{\underline M}$. Then $\widetilde S$ is flat and unramified over $S$. Therefore $e(S/R)=e(\widetilde S/D)$ by lemma 2.3. From lemma 3.6, we get $e(\widetilde S/D)=\sum_{\underline p}L(\widetilde S_{\underline p})e(\widetilde S/\underline p\mid D)$ where $\underline p$ ranges through the minimal primes in $\widetilde S$ such that $\dim\widetilde S/\underline p=\dim\widetilde S$. Now if we let $\underline p_0=\operatorname{Ker}(D\otimes_R S\longrightarrow DS)$, then $\underline p_0\cap R=0$ and thus we are in a position to apply lemma 4.5. Thus we have that $e(\widetilde S/D)=\sum_{\underline p}L(\widetilde S_{\underline p})e(\widetilde S/\underline p\mid D)$ where $\underline p\cap R=0$. We claim that the local $D$-algebras $\widetilde S/\underline p$ are all finite $D$-module. For, since there is one-one correspondence between the prime ideals $\underline p$ in $D\otimes_R S$ with $\underline p\cap R=0$ and the prime ideals in $(D\otimes_R S)\otimes_R k=L'_D\otimes_k K$ (where $L'_D$ is the quotient field of $D$), we see that $\widetilde S/\underline p$ is isomorphic to (and hence we shall identify) a localization of $\sigma(S)D$ as $D$-algebras, for some conjugation $\sigma$ of $S$ over $R$. (The reader may observe that $\sigma(S)D$ corresponding to $\widetilde S/\underline p$ has, by the definition of $\widetilde S=(D\otimes_R S)_{\underline M}$, a maximal ideal which dominates both the maximal ideals of $S^\sigma$ and $D$). The quotient field of $\widetilde S/\underline p$ is $L'_D\sigma(K)$ and $L'_D\sigma(K')$ is the separable closure of $L'_D$ in $L'_D\sigma(K)$. Since $\widetilde S/\underline p$ is a local $D$-algebra, we can find a local $D$-algebra $A$ with quotient field $L'_D\sigma(K')$ such that $D\subset A\subset\widetilde S/\underline p$. Since $L'_D\sigma(K)$ is purely inseparable over $L'_D\sigma(K')$, every integral extension of a local ring $A$ contained in $L'_D\sigma(K)$ is local and consequently local $A$-algebra $\widetilde S/\underline p$ is a finite $A$-module. On the other hand, since the integral closure of $D$ in $L'$ is a finite local $D$-algebra (proposition 1.1), every local $D$-algebra contained in $L'$ is a finite local $D$-algebra. In particular, the local $D$-algebra $A$, being contained in $L'_D\sigma(K')\subset L$, is a finite local $D$-algebra. Therefore we conclude that $\widetilde S/\underline p$ is a finite local $D$-algebra for all $\underline p$ with $\underline p\cap R=0$. Furthermore the residue field of $S/\underline p$, being isomorphic to $\overline S$, is a simple extension of $\overline D=\overline R$ by the hypothesis. It
\mspage{104}
